\renewcommand{\EntradaBibtex}{2026_RecorridoUPV_mas_Jaguar}

\begin{frame}{\citetitle{\EntradaBibtex}$^*$ (1)}
\begin{block}{Problema y objetivo} 
El proyecto evoluciona un simulador de recorridos virtuales del campus de la Universidad Politécnica de Victoria para iOS (anteriormente solo para dispositivos Android).
\end{block} 

Esta tercera iteración busca resolver dos limitaciones:
\begin{itemize}
\item El campus era un espacio sin personajes o vida simulada.
\item El recorrido solamente podía experimentarse en primera persona.Programación y resolución de problemas.
\end{itemize}
Se incorporan:
\begin{itemize}
  \item NPC (Non-Player Character) jaguar autónomo que recorre el campus.
  \item Cámara en tercera persona para visualizar al jugador.
\end{itemize}
\footfullcite*{\EntradaBibtex}
\end{frame}

\begin{frame}{\citetitle{\EntradaBibtex} (2)}
\begin{block}{Arquitectura y componentes} 

Componentes de la arquitectura:
\begin{itemize}
\item UPVScene: controlador principal de la escena.
\item ConstructorEdificio: geometría y sistema de colisiones.
\item Campus y Ambiente: elementos del mundo exterior.
\item Módulos adicionales:
\begin{itemize}
\item Jaguar. Modelo 3D, pelaje procedural y avatar animado
\item JaguarNPC.	Navegación autónoma
\item UPVScene.	Cámara en tercera persona y coordinación
\item ContentView.	Controles y posición del NPC en minimapa
\end{itemize}
\end{itemize}

\end{block} 

\end{frame}

